\documentclass[../../main.tex]{subfiles}
\begin{document}
% SOURCE p=203; printed=188; object=ch06-dual-title; region=upper
\section{$L^p$ 的对偶}\label{section:ch06-lp-dual}
% SOURCE p=203; printed=188; object=ch06-dual-pairing-introduction; region=upper-middle
设 $p,q$ 为共轭指数。Hölder 不等式表明，每个 $g\in L^q$ 都通过
\begin{align*}
\phi_g(f)=\int fg
\end{align*}
在 $L^p$ 上定义一个有界线性泛函，其算子范数至多为 $\|g\|_q$。（若 $p=2$，并把 $L^2$ 看成 Hilbert 空间，那么定义 $\phi_g(f)=\int f\overline g$ 更为合适。同一约定也可用于 $p\ne2$，不会实质改变下面的结果。）事实上，映射 $g\mapsto\phi_g$ 几乎总是从 $L^q$ 到 $(L^p)^*$ 的等距映射。
% SUPPLEMENT: id=v14-ch06-pairing-convention; parent=ch06-dual-pairing-introduction; kind=proof-detail+motivation+hypotheses; origin=independent-derivation
\begin{remark}
\textbf{配对的动机与约定（重构新增）。}积分把两个函数变为标量，Hölder 不等式正好提供连续性估计 $|\phi_g(f)|\leqslant\|f\|_p\|g\|_q$。任一代表在可测零集上的改变不影响乘积积分，所以配对在等价类上良定义。这里采用不取共轭的双线性配对，$g\mapsto\phi_g$ 是复线性映射；改成 Hilbert 约定后，它是共轭线性映射，两者由 $g\mapsto\overline g$ 联系。“几乎总是”不是数学假设：以下命题和定理将分别列出所需的测度及指数条件。
\end{remark}
% END SUPPLEMENT: id=v14-ch06-pairing-convention
% SOURCE p=203; printed=188; object=ch06-dual-pairing-isometry; number=6.13; region=middle
\begin{proposition}\label{proposition:ch06-dual-pairing-isometry}
设 $p,q$ 为共轭指数，且 $1\leqslant q<\infty$。若 $g\in L^q$，则
\begin{align*}
\|g\|_q=\|\phi_g\|=\sup\left\{\left|\int fg\right|:\|f\|_p=1\right\}.
\end{align*}
若 $\mu$ 半有限，则此结果对 $q=\infty$ 也成立。
\end{proposition}
% SOURCE p=203; printed=188; object=ch06-dual-pairing-isometry-proof; region=middle-bottom
\begin{proof}
Hölder 不等式给出 $\|\phi_g\|\leqslant\|g\|_q$，而当 $g=0$ 几乎处处时，等号显然。若 $g\ne0$ 且 $q<\infty$，令
\begin{align*}
f=\frac{|g|^{q-1}\overline{\operatorname{sgn}g}}{\|g\|_q^{q-1}}.
\end{align*}
于是
\begin{align*}
\|f\|_p^p=\frac{\int|g|^{(q-1)p}}{\|g\|_q^{(q-1)p}}
=\frac{\int|g|^q}{\int|g|^q}=1,
\end{align*}
所以
\begin{align*}
\|\phi_g\|\geqslant\int fg=\frac{\int|g|^q}{\|g\|_q^{q-1}}=\|g\|_q.
\end{align*}
（若 $q=1$，则 $f=\overline{\operatorname{sgn}g}$，$\|f\|_\infty=1$，且 $\int fg=\|g\|_1$。）若 $q=\infty$，给定 $\varepsilon>0$，令 $A=\{x:|g(x)|>\|g\|_\infty-\varepsilon\}$。则 $\mu(A)>0$，所以当 $\mu$ 半有限时，存在 $B\subset A$，使 $0<\mu(B)<\infty$。令 $f=\mu(B)^{-1}\chi_B\overline{\operatorname{sgn}g}$，则 $\|f\|_1=1$，因而
\begin{align*}
\|\phi_g\|\geqslant\int fg=\frac{1}{\mu(B)}\int_B|g|\geqslant\|g\|_\infty-\varepsilon.
\end{align*}
由 $\varepsilon$ 的任意性，$\|\phi_g\|=\|g\|_\infty$。
\end{proof}
% SUPPLEMENT: id=v14-ch06-pairing-isometry-details; parent=ch06-dual-pairing-isometry; kind=proof-detail+motivation+hypotheses; origin=independent-derivation
\begin{remark}
\textbf{端点、零空间与逐条件说明（重构新增）。}原证明的 $p$ 次幂计算用于 $1<q<\infty$；$q=1,p=\infty$ 必须用随后括号中的独立计算，不能代入 $(q-1)p$。约定 $\operatorname{sgn}0=0$；零点处直接令测试函数为零，避免 $0^0$。在 $q=\infty$ 的非零情形只需取 $0<\varepsilon<\|g\|_\infty$，这样 $B$ 上 $g\ne0$，测试函数才确有 $L^1$ 范数一。令这样的 $\varepsilon\downarrow0$ 已足够。

若 $L^p=\{0\}$，单位球面为空；本节单位球面公式的非负上确界按 $[0,\infty]$ 中的约定取零。等价且无需空集约定的写法是 $\|\phi_g\|=\sup_{\|f\|_p\leqslant1}|\int fg|$。非零空间中，对任意非零 $f$ 作归一化即可证明两个公式相同；零空间中闭球只含零。在命题条件下，上述非零 $g$ 的测试构造也说明 $L^p=\{0\}$ 必推出 $g=0$。

相位因子把 $fg$ 变成非负数，使积分不发生抵消；有限 $q$ 的幂次由 $(q-1)p=q$ 决定。无穷端点的峰值未必取得，所以用近峰集合。半有限性保证这个正测度集合含正有限测度部分，供 $L^1$ 测试；有限 $q$ 已由 $g$ 自身提供可积测试，无需该假设。若单点空间的非空集测度为无穷，$L^1=\{0\}$，而 $g=1$ 的无穷范数为一，故无穷端点确不能无条件删除半有限性。
\end{remark}
% END SUPPLEMENT: id=v14-ch06-pairing-isometry-details
% SOURCE p=204; printed=189; object=ch06-dual-test-introduction; region=top
反过来，若 $f\mapsto\int fg$ 是 $L^p$ 上的有界线性泛函，那么几乎所有情形下都有 $g\in L^q$。事实上，有如下更强的结果。
% SOURCE p=204; printed=189; object=ch06-finite-support-simple-space; region=upper
\begin{definition}[有限测度支集的简单函数空间]\label{definition:ch06-finite-support-simple-space}
记 $\Sigma$ 为所有在某个有限测度集合之外为零的简单函数所成的空间。
\end{definition}
% SUPPLEMENT: id=v14-ch06-test-space-motivation; parent=ch06-finite-support-simple-space; kind=motivation+hypotheses; origin=independent-derivation
\begin{remark}
\textbf{动机与条件说明（重构新增）。}简单函数使积分可用有限个集合计算；有限测度支集使它同时属于所有 $L^p$，$1\leqslant p\leqslant\infty$。两支集的有限并仍有有限测度，故 $\Sigma$ 对线性运算封闭。这里的“支集”仅指函数在该集合外为零，不涉及拓扑闭包，也不假设 $X$ 有拓扑。
\end{remark}
% END SUPPLEMENT: id=v14-ch06-test-space-motivation
% SOURCE p=204; printed=189; object=ch06-dual-test-criterion; number=6.14; region=upper
\begin{theorem}\label{theorem:ch06-dual-test-criterion}
设 $p,q$ 为共轭指数。设 $g$ 是 $X$ 上的可测函数，对每个 $f\in\Sigma$ 都有 $fg\in L^1$，且
\begin{align*}
M_q(g)=\sup\left\{\left|\int fg\right|:f\in\Sigma\text{ 且 }\|f\|_p=1\right\}<\infty.
\end{align*}
再假设 $S_g=\{x:g(x)\ne0\}$ 为 $\sigma$ 有限集，或者 $\mu$ 半有限。则 $g\in L^q$，且 $M_q(g)=\|g\|_q$。
\end{theorem}
% SOURCE p=204; printed=189; object=ch06-dual-test-criterion-proof; region=middle-bottom
\begin{proof}
首先指出，若 $f$ 为有界可测函数，在有限测度集 $E$ 外为零，且 $\|f\|_p=1$，则 $|\int fg|\leqslant M_q(g)$。事实上，简单逼近定理给出简单函数列 $f_n$，使 $|f_n|\leqslant|f|$（特别地，$f_n$ 在 $E$ 外为零），且 $f_n\to f$ 几乎处处。因为 $|f_n|\leqslant\|f\|_\infty\chi_E$ 且 $\chi_Eg\in L^1$，控制收敛定理给出
\begin{align*}
\left|\int fg\right|=\lim_n\left|\int f_ng\right|\leqslant M_q(g).
\end{align*}
现在设 $q<\infty$。可假设 $S_g$ 为 $\sigma$ 有限集，因为当 $\mu$ 半有限时此条件自动成立，见\cref{exercise:ch06-test-sigma-support}。取有限测度集的递增列 $E_n$，使 $S_g=\bigcup_{n=1}^\infty E_n$。取简单函数列 $\varphi_n$，使 $\varphi_n\to g$ 逐点成立，且 $|\varphi_n|\leqslant|g|$，令 $g_n=\varphi_n\chi_{E_n}$。则 $g_n\to g$ 逐点成立，$|g_n|\leqslant|g|$，且 $g_n$ 在 $E_n$ 外为零。令
\begin{align*}
f_n=\frac{|g_n|^{q-1}\overline{\operatorname{sgn}g}}{\|g_n\|_q^{q-1}}.
\end{align*}
与\cref{proposition:ch06-dual-pairing-isometry}的证明一样，有 $\|f_n\|_p=1$；由 Fatou 引理，
\begin{align*}
\|g\|_q&\leqslant\liminf_n\|g_n\|_q
=\liminf_n\int|f_ng_n|\\
&\leqslant\liminf_n\int|f_ng|
=\liminf_n\int f_ng\leqslant M_q(g).
\end{align*}
最后一个估计使用了证明开头的观察。另一方面，Hölder 不等式给 $M_q(g)\leqslant\|g\|_q$，所以有限 $q$ 的情形得证。

现在设 $q=\infty$。给定 $\varepsilon>0$，令 $A=\{x:|g(x)|\geqslant M_\infty(g)+\varepsilon\}$。若 $\mu(A)>0$，便可取 $B\subset A$，使 $0<\mu(B)<\infty$，这或者由 $\mu$ 半有限，或者由 $A\subset S_g$ 得到。令 $f=\mu(B)^{-1}\chi_B\overline{\operatorname{sgn}g}$，则 $\|f\|_1=1$，且
\begin{align*}
\int fg=\frac{1}{\mu(B)}\int_B|g|\geqslant M_\infty(g)+\varepsilon.
\end{align*}
这与开头的观察矛盾。故 $\|g\|_\infty\leqslant M_\infty(g)$，反向不等式显然。
\end{proof}
% SUPPLEMENT: id=v14-ch06-dual-test-proof-details; parent=ch06-dual-test-criterion; kind=proof-detail+motivation+hypotheses; origin=independent-derivation
\begin{proof}
\textbf{归一化、相位及端点补写（重构新增）。}令 $M=M_q(g)$，空单位球面按上一注记解释。由齐次性，对每个 $u\in\Sigma$ 都有 $|\int ug|\leqslant M\|u\|_p$：范数正时归一化，范数零时 $u=0$ 几乎处处，积分也为零。因此原证的 $f_n$ 不必本身具有单位范数。相同的简单逼近和控制收敛还给出
\begin{align*}
\left|\int ug\right|\leqslant M\|u\|_p
\quad\text{对每个有界、有限测度支集的可测 }u.
\end{align*}
这里 $|u_ng|\leqslant C|g|\chi_E\in L^1$；由 $|u_n|\leqslant|u|$ 有 $\|u_n\|_p\leqslant\|u\|_p$，包括 $p=\infty$。

在原证第二段，若 $\|g_n\|_q=0$，令 $f_n=0$，便有 $\|g_n\|_q=\int|f_ng_n|=0$；不作除以零的操作。若 $1<q<\infty$ 且范数正，原公式在 $g_n=0$ 处取零，得到有界、支于 $E_n$ 的可测函数，且 $\|f_n\|_p=1$。其相位必须取自 $g$：在 $g_n\ne0$ 处 $g\ne0$，所以 $f_ng\geqslant0$，同时 $\int|f_ng_n|=\|g_n\|_q$。若 $q=1$ 且 $\|g_n\|_1>0$，改写为
\begin{align*}
f_n=\chi_{\{g_n\ne0\}}\overline{\operatorname{sgn}g}.
\end{align*}
则 $\|f_n\|_\infty=1$，$f_ng\geqslant0$，$\int|f_ng_n|=\int|g_n|$，原 Fatou 链仍成立。支集示性函数不可省，否则在 $q=1$ 时可能失去有限测度支集，亦会引入 $0^0$。

半有限情形所用的\cref{exercise:ch06-test-sigma-support}只使用本补证第一段的测试估计，独立证明各正水平集有限，并不使用本定理的 $g\in L^q$ 结论，所以没有循环。在无穷端点，若 $S_g=\bigcup_k E_k$ 为有限测度集的可数并，则正测度 $A$ 必有一个 $A\cap E_k$ 正测度，从而可选 $B$。每个 $\{|g|\geqslant M+1/k\}$ 为零集，其可数并包含 $\{|g|>M\}$；因此确有 $\|g\|_\infty\leqslant M$。
\end{proof}
\begin{remark}
\textbf{动机与逐条件说明（重构新增）。}无需预先知道 $g$ 的全局可积性，只要它对所有有限简单测试给出一致有界的积分，就能恢复其 $L^q$ 大小。$g$ 可测保证水平集和相位可测；$fg\in L^1$ 保证局部积分及控制收敛合法；$M<\infty$ 提供统一界，单个测试各自有限并不足够。共轭关系使归一化测试的指数匹配。两个可选测度条件都保证正测度的有关区域能被有限测度测试看到；若二者均无，单点无穷测度与 $g=1$ 给 $M_q(g)=0$，却有 $\|g\|_q=\infty$（有限 $q$）或 $1$（无穷 $q$）。$S_g$ 的 $\sigma$ 有限性只要求 $g$ 非零的区域，不要求整个 $X$。即使 $q=1,p=\infty$，$\Sigma$ 未必在 $L^\infty$ 中稠密，上面的直接测试论证仍然适用。
\end{remark}
% END SUPPLEMENT: id=v14-ch06-dual-test-proof-details
% SOURCE p=204; printed=189; object=ch06-duality-surjection-introduction; region=bottom
对 $(L^p)^*$ 的描述中，最后也是最深刻的一步，是映射 $g\mapsto\phi_g$ 几乎总为满射。
% SOURCE p=205; printed=190; object=ch06-lp-duality; number=6.15; region=top
\begin{theorem}\label{theorem:ch06-lp-duality}
设 $p,q$ 为共轭指数。若 $1<p<\infty$，则对每个 $\phi\in(L^p)^*$，存在 $g\in L^q$，使 $\phi(f)=\int fg$ 对所有 $f\in L^p$ 成立；因此 $L^q$ 与 $(L^p)^*$ 等距同构。若 $\mu$ 为 $\sigma$ 有限测度，则同样的结论对 $p=1$ 成立。
\end{theorem}
% SOURCE p=205; printed=190; object=ch06-lp-duality-proof; region=upper-bottom
\begin{proof}
首先设 $\mu$ 有限，于是所有简单函数都属于 $L^p$。给定 $\phi\in(L^p)^*$，对可测集 $E$ 令 $\nu(E)=\phi(\chi_E)$。若 $E_j$ 两两不交，$E=\bigcup_{j=1}^\infty E_j$，则 $\chi_E=\sum_{j=1}^\infty\chi_{E_j}$，且此级数在 $L^p$ 范数下收敛，因为
\begin{align*}
\left\|\chi_E-\sum_{j=1}^n\chi_{E_j}\right\|_p
=\left\|\sum_{j=n+1}^\infty\chi_{E_j}\right\|_p
=\mu\left(\bigcup_{j=n+1}^\infty E_j\right)^{\frac1p}\longrightarrow0.
\end{align*}
正是在此需要 $p<\infty$。由 $\phi$ 线性且连续，
\begin{align*}
\nu(E)=\sum_{j=1}^\infty\phi(\chi_{E_j})=\sum_{j=1}^\infty\nu(E_j),
\end{align*}
所以 $\nu$ 为复测度。若 $\mu(E)=0$，则 $\chi_E$ 是 $L^p$ 中的零元素，从而 $\nu(E)=0$，即 $\nu\ll\mu$。由 Radon–Nikodym 定理，存在 $g\in L^1(\mu)$，使对所有 $E$ 有 $\phi(\chi_E)=\nu(E)=\int_Eg\,\mathrm d\mu$，从而对所有简单函数 $f$ 有 $\phi(f)=\int fg\,\mathrm d\mu$。又有 $|\int fg|\leqslant\|\phi\|\|f\|_p$，故\cref{theorem:ch06-dual-test-criterion}给出 $g\in L^q$。一旦知道这一点，由\cref{proposition:ch06-simple-density}可知 $\phi(f)=\int fg$ 对所有 $f\in L^p$ 成立。

现在设 $\mu$ 为 $\sigma$ 有限测度。取递增集合列 $E_n$，使 $0<\mu(E_n)<\infty$ 且 $X=\bigcup_{n=1}^\infty E_n$。把 $L^p(E_n)$ 和 $L^q(E_n)$ 分别视为 $L^p(X)$ 和 $L^q(X)$ 中在 $E_n$ 外为零的函数构成的子空间。由刚才的论证，对每个 $n$，存在 $g_n\in L^q(E_n)$，使对所有 $f\in L^p(E_n)$ 有 $\phi(f)=\int fg_n$，且
\begin{align*}
\|g_n\|_q=\|\phi|_{L^p(E_n)}\|\leqslant\|\phi\|.
\end{align*}
$g_n$ 在零集上改值以外是唯一的，所以当 $n<m$ 时，$g_n=g_m$ 在 $E_n$ 上几乎处处成立。可在 $E_n$ 上令 $g=g_n$，从而在 $X$ 上几乎处处定义 $g$。由单调收敛定理，$\|g\|_q=\lim_n\|g_n\|_q\leqslant\|\phi\|$，故 $g\in L^q$。若 $f\in L^p$，控制收敛定理给 $f\chi_{E_n}\to f$ 于 $L^p$ 范数，因而
\begin{align*}
\phi(f)=\lim_n\phi(f\chi_{E_n})=\lim_n\int_{E_n}fg=\int fg.
\end{align*}

最后设 $\mu$ 任意而 $p>1$，从而 $q<\infty$。与上面相同，对每个 $\sigma$ 有限集 $E\subset X$，存在几乎处处唯一的 $g_E\in L^q(E)$，使对所有 $f\in L^p(E)$ 有 $\phi(f)=\int fg_E$，且 $\|g_E\|_q\leqslant\|\phi\|$。若 $F$ 为 $\sigma$ 有限集且 $F\supset E$，则在 $E$ 上 $g_F=g_E$ 几乎处处，所以 $\|g_F\|_q\geqslant\|g_E\|_q$。令 $M$ 为所有 $\sigma$ 有限集 $E$ 对应的 $\|g_E\|_q$ 的上确界，注意 $M\leqslant\|\phi\|$。选取 $E_n$ 使 $\|g_{E_n}\|_q\to M$，令 $F=\bigcup_{n=1}^\infty E_n$。则 $F$ 为 $\sigma$ 有限集，且对每个 $n$ 有 $\|g_F\|_q\geqslant\|g_{E_n}\|_q$，因此 $\|g_F\|_q=M$。若 $A$ 是包含 $F$ 的 $\sigma$ 有限集，则
\begin{align*}
\int|g_F|^q+\int|g_{A\setminus F}|^q
=\int|g_A|^q\leqslant M^q=\int|g_F|^q.
\end{align*}
于是 $g_{A\setminus F}=0$，并且 $g_A=g_F$ 几乎处处。这里用到了 $q<\infty$。但若 $f\in L^p$，则 $A=F\cup\{x:f(x)\ne0\}$ 为 $\sigma$ 有限集，故 $\phi(f)=\int fg_A=\int fg_F$。因此可取 $g=g_F$，证明完成。
\end{proof}
% SUPPLEMENT: id=v14-ch06-duality-proof-details; parent=ch06-lp-duality; kind=proof-detail+motivation+hypotheses; origin=independent-derivation
\begin{proof}
\textbf{三阶段构造的细节（重构新增）。}有限测度情形的尾并递减到空集，故由测度从上连续性，其测度趋零。为额外核验 $\nu$ 的有限变差，任取有限可测分割 $A_1,\ldots,A_N$，令 $c_j=\overline{\operatorname{sgn}\nu(A_j)}$，则
\begin{align*}
\sum_j|\nu(A_j)|=\phi\left(\sum_jc_j\chi_{A_j}\right)
\leqslant\|\phi\|\mu(X)^{\frac1p}.
\end{align*}
左侧为非负实数，故中间式取绝对值不变。这提供统一总变差界。所用复测度版本是\cref{theorem:ch03-complex-lebesgue-radon-nikodym}；因 $\nu\ll\mu$，其奇异部分为零。由简单函数逼近 $f_k\to f$ 于 $L^p$，连续性给 $\phi(f_k)\to\phi(f)$，Hölder 估计给 $\int f_kg\to\int fg$，从而扩展到全部 $f$。若两表示元之差为 $h$，则所有积分测试均为零，\cref{proposition:ch06-dual-pairing-isometry}给 $\|h\|_q=0$，证明唯一性和表示的等距性。

若 $\mu(X)=0$，各空间及其对偶均为零，直接取 $g=0$；原证第二段的正测度穷竭针对 $\mu(X)>0$。后者可先选一个正有限测度集，再把它并入每个有限测度穷竭集。选各 $g_n$ 的可测代表，在不交层 $E_n\setminus E_{n-1}$ 上用 $g_n$ 定义 $g$（$E_0=\varnothing$），便得到处处定义的可测函数，且由相容性，它在每个 $E_n$ 上几乎处处等于 $g_n$。有限 $q$ 时，$\int_{E_n}|g|^q\uparrow\int|g|^q$ 给原范数极限；当 $p=1,q=\infty$ 时，应分别证明：各 $E_n$ 上 $|g|\leqslant\|\phi\|$ 几乎处处，删去可数并的零集便得全局同一界；此外 $\|g\|_\infty=\sup_n\|g\chi_{E_n}\|_\infty$。这一端点不直接套用有限幂积分的单调收敛计算。$p<\infty$ 保证 $\int|f|^p\chi_{E_n^c}\to0$，而 $fg\in L^1$ 保证最后的积分极限。

任意测度阶段，所有 $g_E$ 都在 $E$ 外补零；相容性仅在所比较的两个集合之间逐次使用，无需对不可数多个零集取并。$M=0$ 时可取所有 $E_n=\varnothing$；$M>0$ 时按上确界选 $\|g_{E_n}\|_q>M-1/n$。可数并 $F$ 仍为 $\sigma$ 有限集，因而确有达到上确界的 $g_F$。在原式中，各项至多为有限数 $M^q$，所以相减合法，迫使 $\int_{A\setminus F}|g_A|^q=0$。最后，对任意 $f\in L^p$，
\begin{align*}
\{f\ne0\}=\bigcup_{k\geqslant1}\{|f|>1/k\},\qquad
\mu(\{|f|>1/k\})\leqslant k^p\|f\|_p^p<\infty.
\end{align*}
因此该 $f$ 确被某个 $\sigma$ 有限区域覆盖，足以完成表示；并未要求全空间 $\sigma$ 有限。
\end{proof}
\begin{remark}
\textbf{动机与逐条件说明（重构新增）。}连续泛函先在示性函数上产生测度，再由 Radon–Nikodym 导数恢复“积分权重”。线性给有限可加性，连续性与 $p<\infty$ 把它提升为可数可加性；仅代数线性不足以通过极限。$p\geqslant1$ 使这里使用赋范对偶及 Hölder 理论。有限测度使所有示性函数可测试；$\sigma$ 有限性让有限区域上的表示可通过可数拼接扩展。对于 $p>1$，有限的共轭指数 $q$ 又使新增区域的 $q$ 次质量可相加并强迫其为零，因此最终不需任何半有限或 $\sigma$ 有限假设。$p=1$ 时的 $q=\infty$ 没有这种可加质量，原论证不能删除 $\sigma$ 有限性；下面的例子表明确有反例。这些条件是所列统一结论的充分条件，后面的可分解测度情形还会给出更广的端点版本。
\end{remark}
% END SUPPLEMENT: id=v14-ch06-duality-proof-details
% SOURCE p=206; printed=191; object=ch06-lp-reflexive; number=6.16; region=top
\begin{corollary}\label{corollary:ch06-lp-reflexive}
若 $1<p<\infty$，则 $L^p$ 是自反的。
\end{corollary}
% SUPPLEMENT: id=v14-ch06-reflexive-full-proof; parent=ch06-lp-reflexive; kind=omitted-proof+motivation+hypotheses; origin=independent-derivation
\begin{proof}
\textbf{典范映射的满射性（重构新增）。}令 $q$ 与 $p$ 共轭，并记等距满射 $T:L^q\to(L^p)^*$ 为 $Tg=\phi_g$。任取 $\Psi\in(L^p)^{**}$，则 $\Psi\circ T\in(L^q)^*$。因 $1<q<\infty$，把对偶表示定理用于 $L^q$，得某个 $f\in L^p$，使
\begin{align*}
\Psi(\phi_g)=(\Psi\circ T)(g)=\int gf=\phi_g(f)
\quad(g\in L^q).
\end{align*}
典范嵌入 $J:L^p\to(L^p)^{**}$ 定义为 $Jf(\phi)=\phi(f)$。$T$ 满射表明上式遍及每个 $\phi\in(L^p)^*$，所以 $\Psi=Jf$。因此 $J$ 满射，正是自反性。
\end{proof}
\begin{remark}
\textbf{动机与条件说明（重构新增）。}测量所有连续线性测量后得到的对象，仍来自一个原空间函数。仅有 $L^p$ 与双对偶某个抽象同构并不能证明自反，必须验证上面的典范映射。严格内部指数让 $p,q$ 两次都适用无测度限制的对偶定理；端点不能这样互换，因为 $(L^\infty)^*$ 通常大于 $L^1$。本结论没有要求测度有限或 $\sigma$ 有限，零空间也包含在证明内。
\end{remark}
% END SUPPLEMENT: id=v14-ch06-reflexive-full-proof
% SOURCE p=206; printed=191; object=ch06-endpoint-dual-introduction; region=upper
最后讨论例外情形 $p=1$ 和 $p=\infty$。对任意测度 $\mu$，对应 $g\mapsto\phi_g$ 都把 $L^\infty$ 映入 $(L^1)^*$，但一般既非单射，也非满射。
% SOURCE p=206; printed=191; object=ch06-nonsemifinite-pairing-example; region=upper
\begin{example}\label{example:ch06-nonsemifinite-pairing}
当 $\mu$ 不是半有限测度时，$L^\infty\to(L^1)^*$ 的上述映射不是单射。
\end{example}
\begin{proof}
若 $E\subset X$ 为无限测度集，却不包含正有限测度的子集，且 $f\in L^1$，则 $\{x:f(x)\ne0\}$ 为 $\sigma$ 有限集，因而与 $E$ 相交于零集。因此 $\phi_{\chi_E}=0$，尽管 $\chi_E$ 在 $L^\infty$ 中不是零元素。
\end{proof}
% SUPPLEMENT: id=v14-ch06-nonsemifinite-details; parent=ch06-nonsemifinite-pairing-example; kind=proof-detail+motivation+hypotheses; origin=independent-derivation
\begin{remark}
\textbf{交集为何为零及反例动机（重构新增）。}由半有限性的否定，存在上述可测 $E$。将 $\{f\ne0\}$ 写成有限测度水平集的可数并，每个水平集与 $E$ 的交仍有限测度，又不可能有正有限测度，所以各交均为零集；可数并仍零。因此 $f\chi_E=0$ 几乎处处，积分为零。另一方面 $\mu(E)=\infty>0$，所以 $\|\chi_E\|_\infty=1$。这个区域对所有 $L^1$ 函数不可见，却被通常的 $L^\infty$ 等价关系保留，正是失去单射性的原因。
\end{remark}
% END SUPPLEMENT: id=v14-ch06-nonsemifinite-details
% SOURCE p=206; printed=191; object=ch06-endpoint-injection-remedy; region=upper-middle
这个问题可以通过重新定义 $L^\infty$ 来补救，见\cref{exercise:ch06-linfty-semifinite-23,exercise:ch06-linfty-semifinite-24,exercise:ch06-linfty-semifinite-25}。满射性失败更为微妙，下面用例子说明；另见\cref{exercise:ch06-linfty-dual-cluster}。
% SOURCE p=206; printed=191; object=ch06-countable-cocountable-dual-example; region=middle
\Needspace{6\baselineskip}
\begin{example}\label{example:ch06-countable-cocountable-dual}
设 $X$ 为不可数集，$\mu$ 是 $(X,\mathscr P(X))$ 上的计数测度，$\mathscr M$ 为可数或余可数集构成的 $\sigma$ 代数，$\mu_0$ 是 $\mu$ 在 $\mathscr M$ 上的限制。每个 $f\in L^1(\mu)$ 都在一个可数集外为零，从而 $L^1(\mu)=L^1(\mu_0)$。另一方面，$L^\infty(\mu)$ 由 $X$ 上所有有界函数组成，而 $L^\infty(\mu_0)$ 由那些在某个可数集外为常数的有界函数组成。因此，$L^1(\mu_0)$ 的对偶是 $L^\infty(\mu)$，而不是较小的空间 $L^\infty(\mu_0)$。
\end{example}
% SUPPLEMENT: id=v14-ch06-countable-cocountable-proof; parent=ch06-countable-cocountable-dual-example; kind=omitted-proof+motivation+hypotheses; origin=independent-derivation
\begin{proof}
\textbf{两个函数空间及对偶的完整核验（重构新增）。}对计数测度，可积函数的每个水平集 $\{|f|>1/n\}$ 都有限，因此其非零集可数。任何可数支集函数对 $\mathscr M$ 可测：开集原像若不含零值，则为可数集；若含零值，其补集为可数集。其积分仍为同一绝对收敛的坐标和，故两个 $L^1$ 等距相同。

若 $h$ 是有界 $\mathscr M$ 可测函数，可取有限值简单函数 $s_n$ 一致逼近 $h$。每个 $s_n$ 的有限可测分割中恰有一个余可数块：不可能全可数，因为 $X$ 不可数，也不可能有两个不交的余可数块。因此存在可数集 $C_n$，使 $s_n$ 在 $C_n^c$ 上为常数 $c_n$。令 $C=\bigcup_n C_n$；在非空的 $X\setminus C$ 上固定一点可知 $c_n$ 收敛，并且 $h$ 在整个 $X\setminus C$ 上等于同一极限。反向，任何在可数集外恒定的函数都对 $\mathscr M$ 可测。计数测度只有空零集，所以本质有界就是处处有界。

最后，直接计算 $\ell^1(X)^*$，以免对非 $\sigma$ 有限的 $\mu_0$ 误用端点对偶定理。若 $\phi\in\ell^1(X)^*$，令 $b(x)=\phi(\chi_{\{x\}})$，则 $|b(x)|\leqslant\|\phi\|$。有限支集函数由线性得 $\phi(f)=\sum_xf(x)b(x)$；每个 $\ell^1$ 函数可由其可数支集上的有限截断作范数逼近，所以同一公式对所有 $f$ 成立。反向，任意有界 $b$ 的该和定义有界泛函，其范数恰为 $\sup_x|b(x)|$，下界由单坐标测试给出。选 $A\subset X$，使 $A$ 与 $A^c$ 都不可数，则 $b=\chi_A$ 有界但不对 $\mathscr M$ 可测，故其泛函不能由任何 $L^\infty(\mu_0)$ 元素表示；逐个测试单点便会强迫表示元等于 $b$。
\end{proof}
\begin{remark}
\textbf{动机与条件说明（重构新增）。}可积函数只能探测可数多个坐标，所以缩小可测集合族没有改变 $L^1$；对偶却允许给所有坐标任意有界权重，超出较小 $\sigma$ 代数的可测函数范围。不可数性保证可数与余可数限制确实非平凡；在通常的选择公理下，不可数集可分成两个不可数部分，给出上面的 $A$。计数测度 $\mu_0$ 仍半有限，因为任何非空可测集包含测度一的单点，却不 $\sigma$ 有限：有限测度集只能是有限集，其可数并不能覆盖 $X$。因此半有限性足以保证单射，却不足以保证通常 $L^\infty$ 表示的满射性。
\end{remark}
% END SUPPLEMENT: id=v14-ch06-countable-cocountable-proof
% SOURCE p=206; printed=191; object=ch06-linfty-dual-introduction; region=middle-lower
至于 $p=\infty$，由\cref{proposition:ch06-dual-pairing-isometry}，映射 $g\mapsto\phi_g$ 总是从 $L^1$ 到 $(L^\infty)^*$ 的等距单射，但几乎从来不是满射。\cref{section:ch06-linfty-dual}将进一步讨论；这里先给出具体例子，另一个例子见\cref{exercise:ch06-linfty-dual-cluster}。
% SOURCE p=206; printed=191; object=ch06-linfty-point-evaluation-example; region=lower
\begin{example}\label{example:ch06-linfty-point-evaluation}
令 $X=[0,1]$，$\mu$ 为 Lebesgue 测度。$f\mapsto f(0)$ 是 $C(X)$ 上的有界线性泛函，把 $C(X)$ 视为 $L^\infty$ 的子空间。由 Hahn–Banach 定理，存在 $\phi\in(L^\infty)^*$，使 $\phi(f)=f(0)$ 对每个 $f\in C(X)$ 成立。此 $\phi$ 不能由对某个 $L^1$ 函数积分来表示。
\end{example}
\begin{proof}
考虑 $f_n\in C(X)$，$f_n(x)=\max(1-nx,0)$。对每个 $n$，$\phi(f_n)=f_n(0)=1$；但对每个 $x>0$，$f_n(x)\to0$，所以由控制收敛定理，对每个 $g\in L^1$ 都有 $\int f_ng\to0$。
\end{proof}
% SUPPLEMENT: id=v14-ch06-point-evaluation-details; parent=ch06-linfty-point-evaluation-example; kind=proof-detail+motivation+hypotheses; origin=independent-derivation
\begin{remark}
\textbf{等价类、延拓及条件说明（重构新增）。}连续函数若几乎处处相等，则它们之差若在某点非零，连续性会使差在一个非空相对开区间上保持非零，而该区间有正 Lebesgue 测度，矛盾。因此 $C([0,1])$ 确可等距嵌入 $L^\infty$，其一致范数等于本质上确界。点值泛函先只在此连续子空间定义，范数为一（常数一取得下界），再由\cref{theorem:ch05-complex-hahn-banach}延拓；不能直接对任意 $L^\infty$ 等价类定义 $f(0)$。若延拓存在积分表示 $\phi=\phi_g$，则 $|f_ng|\leqslant|g|\in L^1$ 且 $f_ng\to0$ 几乎处处，故积分应趋零，与恒等于一矛盾。零测度的端点使点值行为与积分行为分离；完整结论只需某个延拓存在，不主张延拓唯一。
\end{remark}
% END SUPPLEMENT: id=v14-ch06-point-evaluation-details
\subfile{section-02-exercises.tex}
\end{document}
